\section{Semantics}

We now move on to discussing the semantics of \ac{PTS}.
We use an untyped domain model, where expressions are viewed as normal forms programs in the untyped $\lambda$-calculus.
Our normal forms are $\beta$-short and $\eta$-long.
We can represent $\beta$-short terms in the usual way, by separating expressions into neutral and normal as follows:
\[
\begin{array}{llcrcl}
  \mbox{Neutral expressions} & \Ne & \ni & E & ::= & x_i \mid \tmapp E V \\
  \mbox{Normal expressions}  & \Nf & \ni & V & ::= & \tmcst s \mid \tppi {V_1} {V_2} \mid \tmlam V \\
\end{array}
\]
Ensuring that expressions are also $\eta$-long is instead achieved through typing rules.

The \ac{NbE} procedure is performed in two steps:
First, syntactic expressions are evaluated into the domain, removing all $\beta$-redeces in the process.
Second, domain objects are readback into normal expressions, performing any necessary $\eta$-expansions in the process.

Semantic judgments are defined using a \ac{PER} model over domain objects.
We adapt \citet{icfp}'s PER model for \mltt, where the \ac{PER} model is defined using impredicativity instead of induction-recursion.
They used this technique in order to mechanize their system in \rocq, which does not support the more familiar inductive-recursive definition.
In our setting, induction-recursion is simply not applicable since \acp{PTS} can be impredicative, in which case the inductive part of the definition is not well-founded.



\subsection{Domain model}

We now define the domain model.

\[
\begin{array}{llcrcl}
  \mbox{Neutral domain objects} & \Dne & \ni & e       & ::=  & z_j \mid \dapp e v \\
  \mbox{Domain objects}         & \Dom & \ni & m,n,a,b & ::=  & \dreflect a e \mid \dcst s \mid \dpi a B \rho \mid \dlam M \rho  \\
  \mbox{Normal domain objects}  & \Dnf & \ni & v       & ::=  & \dreify a m \\
  \mbox{Environments}           & \Env & \ni & \rho    & ::=  & \denvempty \mid \denvext \rho m \\
\end{array}
\]

Like neutral expressions, neutral domain objects consist of variables and function application.
Here, we represent variables $z_j$ with de Bruijn levels (not indices), which means that we count their position from the left of the context instead of the right.
This ensures that variables have a unique identifier throughout an expression, i.e. the level is never shifted, contrary to indices.

Domain object are similar to normal expressions.
They include type-annotated reflected neutrals $\dreflect a e$, sorts, function spaces, and functions themselves.
Functions and their types are defunctionalized by pairing a syntactic expression with an environment interpreting its free variables.
For function spaces $\dpi a B \rho$, the domain $a$ is actually a domain object, but $B$ relies on a variable of type $a$ that has yet to be interpreted, i.e. whose value is not yet in $\rho$.
Likewise, in the function $\dlam M \rho$, the environment $\rho$ provides an interpretation for all variables of $M$ except the newly bound one.
While it is possible to define a domain without referencing any syntactic object, the defunctionalization approach has the advantage of keeping the domain definition inductive.

Finally, normal domain objects consist only of reified domain objects $\dreify a m$,
  where the type annotation $a$ indicates what $\eta$-expansion needs to be performed.
In this setting, reification is not an actual operation performing $\eta$-expansions, but rather a marker that can later guide an actual computation.


\subsection{Normalization functions}

We now move on to defining the normalization procedure.
Since not all terms of all \acp{PTS} are normalizing, we define the procedure using partial functions.
We present them in relational style, as this better lends itself to mechanization.


\subsubsection{Evaluation}

\begin{figure}
  \raggedright
  \noindent
  $\boxed{\evalexp M \rho m}$ -- $M$ evaluates to $m$ in environment $\rho$.
  \[
  \begin{array}{c}
    \infer{
      \evalexp {x_i} \rho {\rho(i)}
    }{
      % Axiom
    }
    \qquad
    \infer{
      \evalexp {\tmcst s} \rho {\dcst s}
    }{
      % Axiom
    }
    \qquad
    \infer{
      \evalexp {\tppi A B} \rho {\dpi a B \rho}
    }{
      \evalexp A \rho a
    }
    \qquad
    \infer{
      \evalexp {\tmlam M} \rho {\dlam M \rho}
    }{
      % Axiom
    }
    \\[8pt]
    \infer{
      \evalexp {\tmapp M N} \rho {m'}
    }{
      \evalexp M \rho m
      &
      \evalexp N \rho n
      &
      \evalapp m n {m'}
    }
    \qquad
    \infer{
      \evalexp {\tmclo \sigma M} \rho m
    }{
      \evalsub \sigma \rho {\rho'}
      &
      \evalexp M {\rho'} m
    }
  \end{array}
  \]

  \noindent
  $\boxed{\evalsub \sigma \rho {\rho'}}$ -- $\sigma$ evaluates to $\rho'$ in environment $\rho$.
  \[
  \begin{array}{c}
    \infer{
      \evalsub \subempty \rho \denvempty
    }{
      % Axiom
    }
    \qquad
    \infer{
      \evalsub \subid \rho \rho
    }{
      % axiom
    }
    \qquad
    \infer{
      \evalsub \subwk {\rho, m} \rho
    }{
      % axiom
    }
    \\[8pt]
    \infer{
      \evalsub {\subext \sigma M} \rho {\rho', m}
    }{
      \evalsub \sigma \rho {\rho'}
      &
      \evalexp M \rho m
    }
    \qquad
    \infer{
      \evalsub {\subclo \sigma \tau} {\rho_1} {\rho_3}
    }{
      \evalsub \sigma {\rho_1} {\rho_2}
      &
      \evalsub \tau {\rho_2} {\rho_3}
    }
  \end{array}
  \]

  \noindent
  $\boxed{\evalapp {m_1} {m_2} {m_3}}$ -- Applying function $m_1$ to $m_2$ yields $m_3$.
  \[
  \begin{array}{c}
    \infer{
      \evalapp {\dreflect {\dpi a B \rho} e} n {\dreflect b {(\dapp e {\dreify a n})}}
    }{
      \evalexp B {\rho, n} b
    }
    \qquad
    \infer{
      \evalapp {\dlam M \rho} m n
    }{
      \evalexp M {\rho, m} n
    }
  \end{array}
  \]
  \caption{Evaluation relations}
  \label{fig:eval}
\end{figure}


The first step in normalization is evaluating syntacting expressions into domain objects.
Evaluation consists of three functions, defined in \Cref{fig:eval}.
The main function is the evaluation of expressions, denoted $\evalexp M \rho m$.
We cannot evaluate the body of functions $\tmlam M$ since the environment $\rho$ does not provide an interpretation for the newly bound variable.
Instead, we simply perform the closure to obtain the domain object $\dlam M \rho$.
The case of function spaces $\tppi A B$ is similar, except that we first evaluate the domain.

Most of the actual work occurs in the function application case,
  which relies on the auxiliairy function that computes function application within the domain, denoted by $\evalapp m n {m'}$.
Application is defined based on the shape of the function $m$.
If $m$ is a concrete function $\dlam M \rho$, then we now have the interpretation $n$ for the newly bound variable, so we can evaluate the body $M$ in the extended environment.
If $m$ is a neutral object $e$ reflected at a function space $\dpi a B \rho$, then we apply $e$ to the reified argument $n$ and reflect the result at the co-domain type,
  which we can correctly evaluate now that we have an interpretation for the newly bound variable.
  
The last function, $\evalsub \sigma \rho {\rho'}$, evaluates substitutions into environments, and is needed to deal with explicit substitutions.
Intuitively, a substitution is a list of expressions, and evaluation simply evaluates each expression in the list.



\subsubsection{Readback}

\begin{figure}
  \raggedright
  \noindent
  $\boxed{\rbnf v V}$ -- Normal domain object $v$ reads back to normal term $V$
  \[
  \begin{array}{c}
    \infer{
      \rbnf {\dreify {\dreflect {a} {e}} {\dreflect b e'}} E
    }{
      \rbne {e'} E
    }
    \qquad
    \infer{
      \rbnf {\dreify {\dcst s} a} V
    }{
      \rbtyp a V
    }
    \qquad
    \infer{
      \rbnf {\dreify {\dpi a B \rho} m} {\tmlam V}
    }{
      \evalexp B {\rho, \dreflect a {z_i}} b
      &
      \evalapp m {\dreflect a {z_i}} n
      &
      \rbnf[i+1] {\dreify b n} V
    }
  \end{array}
  \]

  \noindent
  $\boxed{\rbne e E}$ -- Neutral domain object $e$ reads back to neutral term $E$.
  \[
  \begin{array}{c}
    \infer{
      \rbne {z_j} {x_{i - j + 1}}
    }{
      % Axiom
    }
    \qquad
    \infer{
      \rbne {\dapp e v} {\tmapp E V}
    }{
      \rbne e E
      &
      \rbnf v V
    }
  \end{array}
  \]

  \noindent
  $\boxed{\rbtyp a V}$ -- Domain type $a$ readbacks to normal type $V$.
  \[
  \begin{array}{c}
    \infer{
      \rbtyp {\dreflect {\tmcst s} e} E
    }{
      \rbne e E
    }
    \qquad
    \infer{
      \rbtyp {\tmcst s} {\tmcst s}
    }{
      % Axiom
    }
    \qquad
    \infer{
      \rbtyp {\dpi a B \rho} {\tppi V {V'}}
    }{
      \rbtyp a V
      &
      \evalexp B {\rho, \dreflect a {z_i}} b
      &
      \rbtyp[i+1] b {V'}
    }
  \end{array}
  \]
  \caption{Readback relations}
  \label{fig:readback}
\end{figure}


Next, we define readback relations that translate domain objects back into $\eta$-expanded syntactic normal forms.
There are three kinds of readbacks: one for normal objects, one for neutral objects, and one types.
Each readback is indexed by a natural number $i$ corresponding to the number of variables in the context.
This index serves two purposes:
First, it is needed by the neutral readback to translate de Bruijn levels back into de Bruijn indices;
Second, it tells what the level of a fresh variable should be.

The main readback is the one for normal domain objects, denoted as $\rbnf v V$.
The normal readback is defined based on the type at which the object is reified.
If the type is a sort $\dcst s$, then we simply invoke the readback specialized in dealing with types.
If it is a function space $\dpi a B \rho$, then we use the domain application relation to apply the function to a fresh variable, readback the result, and wrap it in a lambda.
This is how $\eta$-expansion is performed.

The readback for neutral domain objects, denoted $\rbne e E$, mainly performs the conversion from de Bruijn levels to de Bruijn indices.
It $\Gamma$ is a context of length $i$ and $j < i$, then the $j^{th}$ position from the left is the $(i - j + 1)^{th}$ position from the right.
That is, a variable with de Bruijn level $j$ has de Bruijn index $i - j + 1$.
To readback a function application, we simply use the appropriate readbacks on the function and arguments.

Finally, we have the readback specialized in dealing with types, denoted $\rbtyp a V$.
This readback is straightforward:
The case of a reflected neutral simply invokes the neutral readback;
The case of a sort simply reads it back to itself;
The case of a function space is essentially a congruence on the domain and co-domain.



\subsubsection{Normalization}

Given a derivation of $\wfexp \Gamma M A$, we want to compute the normal form of $M$.
Since $M$ is well-formed in context $\Gamma$, it can contain free variables.
Moreover, the necessary $\eta$-expansions are given by the type $A$.
Consequently, we cannot simply normalize $M$ in isolation, and we must intead take into account the entire information contained in the typing judgment.

To deal with the free variables, we need to define an initial environment for $\Gamma$, which we denote as $\denvinit \Gamma$.
Intuitively, $\denvinit \Gamma$ is an identity environment, that is, it consists only of variables, suitable reflected at the types in $\Gamma$.
Formally, we define initial environments as follows:
\[
\begin{array}{lcll}
%  \denvinit \_                 & : & \Ctx \to \Env \\
  \denvinit \ctxempty          & = & \denvempty \\
  \denvinit {\ctxext \Gamma A} & = & \denvext {\denvinit \Gamma} {\dreflect a {z_{|\Gamma|}}} & \text{where } \evalexp A {\denvinit \Gamma} a \\[8pt]
\end{array}
\]

Now, we can define a normalization relation $\nbeexp \Gamma M A V$, meaning that an expression $M$ of type $A$ in context $\Gamma$ has normal form $V$.
This relation is define simply by composing the evaluation and readback relations, as follows:
\[
\infer{
  \nbeexp \Gamma M A V
}{
  \evalexp A {\denvinit \Gamma} a
  &
  \evalexp M {\denvinit \Gamma} m
  &
  \rbnf[|\Gamma|] {\dreify a m} V
}
\]



\subsubsection{Properties of evaluation relations}

\begin{lemma}[Determinacy]
  All evaluation, readback, and normalization judgments are deterministic.
  Specifically:
  \begin{enumerate}
  \item If $\evalexp M \rho m$ and $\evalexp M \rho {m'}$, then $m = m'$.
  \item If $\evalsub \sigma \rho {\rho_1}$ and $\evalsub \sigma \rho {\rho_2}$, then $\rho_1 = \rho_2$.
  \item If $\evalapp {m_1} {m_2} {m_3}$ and $\evalapp {m_1} {m_2} {m_3'}$, then $m_3 = m_3'$.
  \item If $\rbne e E$ and $\rbne e {E'}$, then $E = E'$.
  \item If $\rbnf v V$ and $\rbnf v {V'}$, then $V = V'$.
  \item If $\rbtyp a V$ and $\rbtyp a {V'}$, then $V = V'$.
  \item If $\nbeexp \Gamma M A V$ and $\nbeexp \Gamma M A {V'}$, then $V = V'$.
  \end{enumerate}
\end{lemma}
\begin{proof}
  By induction on the first derivation.
\end{proof}

This lemma tells us that we can view the judgments as partial functions.
We denote the functional versions by omitting the $\searrow$ symbol and the output.
For example, if $\evalexp M \rho m$, then we write $\evalexpfunc M \rho$ to mean $m$.


\subsection{PER model}

The last things we need to define are the semantic analogues to our syntactic judgments.
This is achieved using a \ac{PER} model, i.e. symmetric and transitive relations that represent types.
Given a \ac{PER} $\per A$, we denote the relation of objects as $\perrel m n {\per A}$, and interpret it as meaning that $m$ and $n$ are equal at type $\per A$.
We write $m \in \per A$ as shorthand for $\perrel m m {\per A}$, and interpret this as meaning that $m$ has type $\per A$.
While \acp{PER} are not reflexive over the entire domain, it follows from symmetry and transitivity that if $\perrel m n {\per A}$, then $m \in \per A$.

The key idea behind defining a \ac{PER} model is that we want to represent the intensional syntactic judgments in an extensional style.
For example, two functions are syntactically equal if their bodies are syntactically equal in an extended context, where the new variable is only equal to itself.
In contrast, two functions are semantically equal if they map any semantically equal inputs to semantically equal outputs.




\subsubsection{PERs for neutral and normal domain objects}

\[
\begin{array}{l@{~\iff~}l}
  \perrel e {e'} \perne  & \forall i \in \Nat.\exists E \in \Ne.\ \rbne e E ~ \text{and} ~ \rbne {e'} E \\
  \perrel v {v'} \pernf  & \forall i \in \Nat.\exists V \in \Nf.\ \rbnf v V ~ \text{and} ~ \rbnf {v'} V \\
  \perrel a {a'} \pertyp & \forall i \in \Nat.\exists V \in \Nf.\ \rbtyp a V ~ \text{and} ~ \rbtyp {a'} V \\[12pt]
  \perrel {\dreflect a e} {\dreflect {a'} {e'}} \perneu & \perrel e {e'} \perne
\end{array}
\]


\subsubsection{PERs for sorts}

The main challenge with defining the \ac{PER} model for a \ac{PTS} is providing a \ac{PER} for sorts.
Due to the extensional nature of the judgment, relating two function spaces requires us to extract a \ac{PER} interpretation of their domains and co-domains,
  so that we can meaningfully speak of related inputs and outputs.
This means that a sort essentially needs to be represented as a \ac{PER} of \acp{PER}, while simultaneously being a relation over domain objects.
A common approach to tackle this problem is to define an inductive relation over domain objects simultaneously with a recursive function mapping domain objects to \acp{PER}.
However, this only works for predicative type systems, while our system should be general enough to deal with impredicativity.

\[
\begin{array}{c}
  \infer{
    \tmcst s = \tmcst s \in \persort[\tmcst s'] \searrow R
  }{
    (\tmcst s : \tmcst s') \in \Ax
    &
    R \iff \persort
  }
  \qquad
  \infer{
    \dreflect {\tmcst s} e = \dreflect {\tmcst s} {e'} \in \persort \searrow R
  }{
    e = e' \in \perne
    &
    R \iff \perneu
  }
  \\[8pt]
  \infer{
    \dpi a B \rho = \dpi {a'} {B'} {\rho'} \in \persort[\tmcst s_3] \searrow R
  }{
    \begin{array}{c}
      (\tmcst s_1, \tmcst s_2, \tmcst s_3) \in \Ru
      \qquad
      a = a' \in \persort[\tmcst s_1] \searrow R_I
      \qquad
      \mhighlight {R_I \text{ is PER}}
      \\
      \forall \DD :: n = n' \in R_I.\ \evalexpfunc B {\denvext \rho n} = \evalexpfunc {B'} {\denvext {\rho'} {n'}} \in \persort[\tmcst s_2] \searrow R_O(\DD)
      \\
      \forall m, m' \in \Dom.\ m = m' \in R \iff (\forall \DD :: n = n' \in R_I.\ \evalappfunc m n = \evalappfunc {m'} {n'} \in R_O(\DD))
    \end{array}
  }
\end{array}
\]

\subsubsection{PERs for contexts}

\[
\begin{array}{c}
  \infer{
    \ctxempty = \ctxempty \in \perctx \searrow R
  }{
    \forall \rho, \rho'. \rho = \rho' \in R \iff \rho = \rho' = \denvempty
  }
  \\[8pt]
  \infer{
    \ctxext \Gamma A = \ctxext \Delta B \in \perctx \searrow R
  }{
    \begin{array}{c}
      \Gamma = \Delta \in \perctx \searrow R'
      \qquad
      \mhighlight{R' \text{ is PER}}
      \\
      \forall \DD :: \rho = \rho' \in R'.\ \evalexpfunc A \rho = \evalexpfunc {A'} {\rho'} \in \persort \searrow R''(\DD)
      \\
      \forall \rho, \rho', m, m'.\ \denvext \rho m = \denvext {\rho'} {m'} \in R \iff (\exists \DD :: \rho = \rho' \in R' ~ \text{and} ~ m = m' \in R''(\DD))
    \end{array}
  }
\end{array}
\]


\subsubsection{Properties of PER model}


